% ============================================================
%  glossary.tex — Inside a Smartphone: The Tiny World in Your Hand
%  Alphabetical glossary. Target: 30–35 terms.
%  Add \emph{(Chapter N)} after each definition.
% ============================================================

\namedchapter{Glossary}
\addcontentsline{toc}{chapter}{Glossary}

\begin{tcolorbox}[enhanced,breakable,colback=SoftPurple!5,colframe=SoftPurple!60,arc=6pt,boxrule=1pt,
  left=10pt,right=10pt,top=10pt,bottom=10pt,title={\faBook\quad All the Words You Learned},
  fonttitle=\bfseries\Large\color{DarkText},before upper=\RaggedRight]

\begin{description}[labelwidth=4cm, leftmargin=4.2cm, style=nextline, itemsep=0.45\baselineskip]

\item[\glossterm{Accelerometer}]
A sensor that measures how fast a phone is speeding up or slowing down.
\emph{(Chapter 7)}

\item[\glossterm{App}]
A programme that runs on a smartphone to do a specific job — like playing music or showing a map.
\emph{(Chapter 8)}

\item[\glossterm{Battery}]
A device that stores chemical energy and converts it into electricity to power the phone.
\emph{(Chapter 6)}

\item[\glossterm{Biometrics}]
Using a unique part of your body — like your fingerprint or face — to unlock a device.
\emph{(Chapter 10)}

\item[\glossterm{Camera sensor}]
The part of a camera that converts light into digital information.
\emph{(Chapter 5)}

\item[\glossterm{Capacitive touchscreen}]
A type of screen that detects the tiny electrical charge in a human fingertip.
\emph{(Chapter 4)}

\item[\glossterm{Charging}]
Restoring energy to a battery by connecting it to a power source.
\emph{(Chapter 6)}

\item[\glossterm{Clock speed}]
How many instructions a processor can carry out every second, measured in gigahertz (GHz).
\emph{(Chapter 2)}

\item[\glossterm{Compass}]
A sensor that detects Earth's magnetic field to determine which direction the phone is facing.
\emph{(Chapter 7)}

\item[\glossterm{GPS}]
Global Positioning System — a network of satellites that helps devices find their location.
\emph{(Chapter 9)}

\item[\glossterm{Gyroscope}]
A sensor that measures the rotation and tilt of a phone.
\emph{(Chapter 7)}

\item[\glossterm{Internal storage}]
The permanent memory inside a phone that keeps photos, apps, and files even when the phone is off.
\emph{(Chapter 3)}

\item[\glossterm{Lens}]
A curved piece of glass or plastic that focuses light onto the camera sensor.
\emph{(Chapter 5)}

\item[\glossterm{Light sensor}]
A sensor that measures the brightness of the surroundings and adjusts screen brightness automatically.
\emph{(Chapter 7)}

\item[\glossterm{Lithium-ion battery}]
The common type of rechargeable battery used in smartphones.
\emph{(Chapter 6)}

\item[\glossterm{Microphone}]
A sensor that converts sound waves into electrical signals so the phone can record or transmit voice.
\emph{(Chapter 7)}

\item[\glossterm{Night mode}]
A camera feature that takes multiple photos quickly and combines them to produce a brighter image in low light.
\emph{(Chapter 5)}

\item[\glossterm{Operating system}]
The main programme that manages all the hardware and software on a phone (e.g., Android, iOS).
\emph{(Chapter 8)}

\item[\glossterm{Pixel}]
The smallest dot of colour on a screen or in a digital image. More pixels means more detail.
\emph{(Chapter 5)}

\item[\glossterm{Processor}]
The brain of the phone — a tiny chip that carries out instructions and controls everything the phone does.
\emph{(Chapter 2)}

\item[\glossterm{RAM}]
Random Access Memory — the phone's working memory, used to hold information the phone is
actively using right now.
\emph{(Chapter 3)}

\item[\glossterm{Satellite}]
A machine that orbits Earth and sends or receives signals to and from devices on the ground.
\emph{(Chapter 9)}

\item[\glossterm{Screen resolution}]
The number of pixels on a screen. Higher resolution means sharper images.
\emph{(Chapter 4)}

\item[\glossterm{Selfie}]
A photo taken of yourself, usually with the front-facing camera.
\emph{(Chapter 5)}

\item[\glossterm{Sensor}]
A tiny device inside a phone that detects changes in the physical world — like movement, light, or direction.
\emph{(Chapter 7)}

\item[\glossterm{Smartphone}]
A mobile phone that is also a computer — it can run apps, connect to the internet, and use cameras and sensors.
\emph{(Chapter 1)}

\item[\glossterm{Storage}]
The space used to save information permanently on a phone.
\emph{(Chapter 3)}

\item[\glossterm{Touchscreen}]
A screen that can detect when and where you touch it, turning your finger into a control.
\emph{(Chapter 4)}

\item[\glossterm{Transistor}]
A microscopic switch inside a processor that can be on or off — billions of them work together to process information.
\emph{(Chapter 2)}

\item[\glossterm{Triangulation}]
Using signals from three or more satellites to calculate an exact position on Earth.
\emph{(Chapter 9)}

\item[\glossterm{Update}]
A new version of an app or the operating system that fixes problems or adds features.
\emph{(Chapter 8)}

\item[\glossterm{Zoom}]
Making part of an image appear larger. Digital zoom crops the image; optical zoom uses moving lenses.
\emph{(Chapter 5)}

\end{description}
\end{tcolorbox}
